\section{Systems of Linear Equations}

\begin{defbox}
Solve $[A]\{X\}=\{B\}$: $n$ equations, $n$ unknowns. \textbf{Gauss elimination} has two steps: (1) \emph{forward elimination} — knock out unknowns until the system is upper triangular; (2) \emph{back substitution} — solve the last (single-unknown) equation, then substitute upward.
\end{defbox}

\subsection{Naive Gauss Elimination}

Multiplier: $m_{i1}=a_{i1}/a_{11}$, used via $R_i\leftarrow R_i-m_{i1}R_1$. General step $k$:
\[
m_{ik} = \frac{a_{ik}^{(k-1)}}{a_{kk}^{(k-1)}}, \qquad a_{ij}^{(k)} = a_{ij}^{(k-1)} - m_{ik}a_{kj}^{(k-1)}, \quad i,j>k
\]
Back substitution: $x_n = b_n^{(n-1)}/a_{nn}^{(n-1)}$, then $x_i = \left(b_i^{(i-1)}-\sum_{j>i}a_{ij}^{(i-1)}x_j\right)/a_{ii}^{(i-1)}$ for $i=n-1,\dots,1$.

\begin{examplebox}
\textbf{Worked example — rocket velocity fit.} Fit $v(t)=a_1t^2+a_2t+a_3$ to $(5,106.8),(8,177.2),(12,279.2)$:
\[
\begin{bmatrix}25&5&1\\64&8&1\\144&12&1\end{bmatrix}\!\begin{bmatrix}a_1\\a_2\\a_3\end{bmatrix}=\begin{bmatrix}106.8\\177.2\\279.2\end{bmatrix}
\]
Step 1 ($m_{21}=2.56,\ m_{31}=5.76$):
\[
\begin{bmatrix}25&5&1&106.8\\0&-4.8&-1.56&-96.208\\0&-16.8&-4.76&-335.968\end{bmatrix}
\]
Step 2 ($m_{32}=3.5$):
\[
\begin{bmatrix}25&5&1&106.8\\0&-4.8&-1.56&-96.208\\0&0&0.7&0.76\end{bmatrix}
\]
Back substitution: $a_3=1.08571$, $a_2=19.6905$, $a_1=0.290472$.
\[
v(t)=0.290472t^2+19.6905t+1.08571
\]
\end{examplebox}

\begin{keypoint}
\textbf{Two failure modes of naive Gauss elimination.}
\begin{enumerate}
\item \textbf{Division by zero.} E.g.\ $a_{11}=0$ from the start — nothing to pivot on, even though the system is perfectly solvable. Sneakier version: a pivot that is nonzero initially can become \emph{exactly zero after an earlier elimination step} (e.g.\ $a_{22}'=0$ appears only after eliminating column 1).
\item \textbf{Round-off error.} Dividing by a very small pivot amplifies tiny rounding differences catastrophically. Worked case: system with exact answer $(1,1,1)$; a pivot of $0.001$ appears at step 2, giving multiplier $-2750$. Rounding the next entry to \textbf{6} vs.\ \textbf{5} significant digits gives final answers $(0.9625,1.05,0.999995)$ (max error $\approx5\%$) vs.\ $(0.625,1.5,0.99995)$ (max error $\approx50\%$) — \emph{one fewer digit of precision, an entirely wrong answer.}
\end{enumerate}
\end{keypoint}

\subsection{Partial Pivoting}

\begin{defbox}
At the start of each elimination step $k$, scan column $k$ for the largest $|a_{ik}|$ ($i\ge k$) and swap that row into the pivot position before eliminating. Fixes \emph{both} failure modes at once: no zero pivots, and much smaller multipliers $\Rightarrow$ far less round-off.
\end{defbox}

\begin{examplebox}
\textbf{Same round-off example, with pivoting.} At step 2, candidate pivots are $|0.001|$ vs.\ $|-2.75|$ — swap rows: pivot becomes $-2.75$, multiplier $m_{32}=0.001/{-2.75}\approx-0.00036$ (\emph{seven million times smaller} than the naive $-2750$). Result with only 5 significant digits: $x_1=x_2=x_3=1$ — \textbf{exact}. ``Same five digits, same system. Only difference: we swapped one row.''
\end{examplebox}

\subsection{Determinant via Elimination}

Three theorems make $\det(A)$ almost free after elimination:
\begin{enumerate}
\item Adding/subtracting a multiple of one row to another leaves $\det$ unchanged.
\item Triangular/diagonal $[A]$: $\det(A)=\prod_i a_{ii}$.
\item Swapping two rows flips the sign of $\det$.
\end{enumerate}

\begin{examplebox}
\textbf{Worked example.} Rocket matrix $\to$ triangular form $\begin{bmatrix}25&5&1\\0&-4.8&-1.56\\0&0&0.7\end{bmatrix}$: $\det(A)=25\times(-4.8)\times0.7=-84.00$ (Theorems 1, 2). With one row swap during pivoting elsewhere: $\det(B)=10\times2.5\times6.002=150.05\Rightarrow\det(A)=-150.05$ (Theorem 3).
\end{examplebox}

\subsection{Gauss--Jordan Elimination}

\begin{defbox}
One step further than Gauss elimination: (i) eliminate the pivot variable from \emph{every} other row (above and below, not just below); (ii) normalize each pivot row so every pivot becomes 1. Result: the coefficient matrix becomes the identity — the RHS \emph{is} the solution, no back substitution needed. Trade-off: strictly more arithmetic than plain Gauss elimination.
\end{defbox}

\begin{examplebox}
\textbf{Worked example.} $3x_1{-}0.1x_2{-}0.2x_3=7.85,\ 0.1x_1{+}7x_2{-}0.3x_3={-}19.3,\ 0.3x_1{-}0.2x_2{+}10x_3=71.4$. After normalizing/eliminating pivot columns 1, 2, 3 in turn:
\[
\begin{bmatrix}1&0&0&3.0000\\0&1&0&-2.5000\\0&0&1&7.0000\end{bmatrix}
\]
Answer read directly off the RHS: $x_1=3,\ x_2=-2.5,\ x_3=7$.
\end{examplebox}

\subsection{LU Decomposition}

\begin{defbox}
$A=LU$: $U$ is exactly the upper-triangular result of naive Gauss forward elimination; $L$ is unit-lower-triangular, built ``for free'' by storing each multiplier $m_{ij}$ at position $(i,j)$ instead of discarding it. Solving $AX=C$ becomes $LZ=C$ (forward substitution, no division — $L$ has 1's on the diagonal) then $UX=Z$ (back substitution, one division per row).
\end{defbox}

\begin{examplebox}
\textbf{Worked example (rocket matrix again).} Step 1 ($m_{21}{=}2.56,m_{31}{=}5.76$), step 2 ($m_{32}{=}3.5$):
\[
U=\begin{bmatrix}25&5&1\\0&-4.8&-1.56\\0&0&0.7\end{bmatrix}, \qquad L=\begin{bmatrix}1&0&0\\2.56&1&0\\5.76&3.5&1\end{bmatrix}
\]
Forward sub.\ with $C=[106.8,177.2,279.2]^\top$: $z_1=106.8,\ z_2=-96.208,\ z_3=0.76$ (identical to Gauss elimination's RHS column — not a coincidence: $Z=L^{-1}C$). Back sub.\ on $UX=Z$ recovers the same $a_1,a_2,a_3$ as before.

\textbf{Finding $A^{-1}$ via LU.} Solve $A\{b_{\cdot j}\}=\{e_j\}$ for each unit column $e_j$ (same $L,U$, only RHS changes):
\[
A^{-1} = \begin{bmatrix}0.04762&-0.08333&0.03571\\-0.9524&1.417&-0.4643\\4.571&-5.000&1.429\end{bmatrix}
\]
\end{examplebox}

\begin{keypoint}
\textbf{Cost comparison (clock-cycle model: $+,-,\times=4T$, $\div=16T$).} For a \emph{single} solve, $C_T|_{LU}=C_T|_{GE}$ exactly — LU buys nothing over Gauss elimination. But for computing $A^{-1}$ (equivalently, $n$ different right-hand sides): LU decomposes once ($\sim n^3$) and reuses it for $n$ cheap ($\sim n^2$) substitutions — total $\sim n^3$; Gauss must redo the full $\sim n^3$ elimination for every RHS — total $\sim n^4$. Ratio $\approx n/4$ (e.g.\ $2501\times$ at $n=10{,}000$). \textbf{``Decompose once, reuse often.''}
\end{keypoint}

\begin{qabox}
\textbf{Q1:} What are the two big steps of Gauss elimination? \\
\textbf{A:} Forward elimination (reduce to upper triangular) followed by back substitution (solve from the last equation upward).

\textbf{Q2:} Why can naive Gauss elimination fail even if no original coefficient is zero? \\
\textbf{A:} A pivot can become exactly zero \emph{after} earlier elimination steps modify the matrix, causing division by zero later even though the original matrix had no zero pivot candidates.

\textbf{Q3:} What is partial pivoting, and what two problems does it fix simultaneously? \\
\textbf{A:} Swap in the row with the largest $|$coefficient$|$ in the current column before eliminating; this avoids zero pivots and greatly reduces round-off error (much smaller multipliers).

\textbf{Q4:} How do you get $\det(A)$ ``for free'' from Gauss elimination? \\
\textbf{A:} Row addition/subtraction doesn't change $\det$; once triangular, $\det(A)=\prod a_{ii}$; each row swap during pivoting flips the sign.

\textbf{Q5:} What two upgrades does Gauss--Jordan add over Gauss elimination, and at what cost? \\
\textbf{A:} Eliminates the pivot variable from every row (not just below) and normalizes every pivot to 1, so the RHS becomes the solution directly (no back substitution) — at the cost of strictly more arithmetic.

\textbf{Q6:} What is the LU decomposition of $A$, and how is it obtained? \\
\textbf{A:} $A=LU$ where $U$ is the upper-triangular result of naive Gauss forward elimination and $L$ is unit-lower-triangular holding the multipliers $m_{ij}$ at position $(i,j)$ — obtained for free during elimination.

\textbf{Q7:} Is LU decomposition cheaper than Gauss elimination for a single solve? \\
\textbf{A:} No — $C_T|_{LU}=C_T|_{GE}$ exactly for one right-hand side.

\textbf{Q8:} When does LU actually pay off? \\
\textbf{A:} With many right-hand sides sharing the same $A$ (e.g.\ computing $A^{-1}$): decompose once ($\sim n^3$) and reuse cheap ($\sim n^2$) substitutions each time, vs.\ Gauss's $\sim n^4$ total — ratio $\approx n/4$ for large $n$.

\textbf{Q9:} Why is division especially expensive in the cost model, and where does this matter? \\
\textbf{A:} Division costs 16 cycles vs.\ 4 for $+,-,\times$ in the model used; this is why forward substitution (unit-diagonal $L$, no division) is cheaper than back substitution (one division per row).
\end{qabox}
